% Осень26!B29; нумерация внутри столбца без пустых строк.
% Источник: лекция 5, §9 (строка 278)
\definitionticket{26}{Интеграл суммируемой функции}
\label{definition:26}

\subsubsection*{Определение 3: суммируемость и интеграл общего вида}

\begin{definition}

\emph{3.} Пусть \(f:X\to\overline{\mathbb R}\) --- измеримая функция произвольного знака, \(f^+=\max(f,0)\), \(f^-=\max(-f,0)\).

\begin{itemize}
\tightlist
\item
  Если \(\displaystyle\int_X f^+\,d\mu<+\infty\) \textbf{и} \(\displaystyle\int_X f^-\,d\mu<+\infty\), функция \(f\) называется \textbf{суммируемой}, и \[\int_X f\,d\mu:=\int_X f^+\,d\mu-\int_X f^-\,d\mu\in\mathbb R.\]
\item
  Если \textbf{ровно одна} из величин \(\int f^+\,d\mu\), \(\int f^-\,d\mu\) равна \(+\infty\), та же формула определяет \(\displaystyle\int_X f\,d\mu=\pm\infty\) (функция \textbf{не} суммируема, но интеграл определён).
\item
  Если обе равны \(+\infty\), формула ничего разумного предложить не может (\(\infty-\infty\)); интеграл не определяется.
\end{itemize}

\end{definition}

\textbf{О терминологии.} Слово «интегрируемая» к моменту создания теории было «затаскано» (использовалось во всяких разных смыслах), поэтому ввели новое слово --- \textbf{суммируемая}: интеграл корректно определён и вычисляется как \textbf{конечное} число. Аналогия с рядами: ряду иногда приписывают сумму \(+\infty\) --- это не мешает считать его расходящимся.

\subsubsection*{Замечания к определению 3}

\begin{itemize}
\item
  \textbf{Совместимость 3 ↔ 2:} если \(f\ge0\), то \(f^+=f\), \(f^-=0\), \(\int f^-\,d\mu=0\), и определение 3 воспроизводит определение 2. Противоречий нет.
\item
  \textbf{Анонс (не доказано, «почему могу»):} \(f\) суммируемая \(\iff\) \(\displaystyle\int_X|f|\,d\mu<+\infty\).

  В одну сторону «почти сразу»: если \(\int|f|<\infty\), то срезки, для которых \(f^\pm\le|f|\), тоже имеют конечные интегралы (нужна линейность/монотонность интеграла для неотрицательных функций). В другую --- понадобится линейность интеграла (складывать интегралы), её ещё нет: «вот сейчас научимся складывать --- и замечание вступит в полную силу».
\end{itemize}

\emph{(Далее перерыв --.)}
